\scanpage{70}
\mspage{1}
\margin{71}
\uline{1}. Soit $R$ un ensemble (ce sera bientôt $\mathbb R$). \deleted{\unsure{Soit C}}
Soit pour \deleted{tout} $n\in\mathbb N$, $C_n\unsure{\in|\subset}R^n$. On suppose

\textcircled{a} $C_n$ stable par $\mathfrak S_n$

\textcircled{b} Soient $p,n\in\mathbb N$, $p\leq n$, considérons
\[
\varphi_{p,n}:R^n\longrightarrow R^p,\qquad
(x_1,--x_n)\longmapsto(x_1,--,x_p)
\]
alors $\varphi_{p,n}(X)\in C_p$ si $X\in C_n$

\textcircled{c} Soient $p_1,p_2,--,p_n\in\mathbb N$, \inserted{$p_*=(p_1,--p_n)$}, considérons
\[
R^n\xrightarrow{\varphi_{p_*}}R^{\overbrace{p_1+\cdots+p_n}^{p}}
\]
\[
x_1,-x_n\longmapsto
(\underbrace{x_1,--x_1}_{p_1\ \mathrm{fois}},--,\underbrace{x_n,--x_n}_{p_n\ \mathrm{fois}})
\]
\deleted{d) $\exists n$ tel que $C_n\ne\varnothing$ $(\Leftrightarrow C_0\ne\varnothing$}
alors $\varphi_{p_*}(X)\in C_p$ si $X\in C_n$

\textcircled{d} Si $X\in C_m$, $Y\in C_n$, alors $X\times Y\in C_{m+n}$

\uline{NB} Il revient au $=$, pour tt ens. fini $I$, de se donner \deleted{\unsure{le}} $C_I\subset\mathcal P(R^I)$, de telle façon que pour tt $\alpha:I\to J$ ($I,J$ ens finis) d'où $\alpha^*:R^J\to R^I$, on ait $\alpha^*(X)\in C_I$ pour $X\in C_J$, et que $X\in C_I$, $Y\in C_J\Rightarrow X\times Y\in C_{I\amalg J}$.

\uline{Prop} Si $X,Y\in C_n$, alors $X\cap Y\in C_n$

\uline{Déf} Soit $X\in C_I$, $Y\in C_{\unsure{J|Y}}$, $f:X\to Y$, d'où $\Gamma_f\subset X\times Y\subset R^{I\amalg J}$. On dit que $f$ est un $C_*$-morphisme si $\Gamma_f\in C_{I\amalg J}$.
\inserted{[NB. ne dépend que du composé \deleted{et non de}
\diagram{\begin{tikzcd}[ampersand replacement=\&]
X\arrow[r]\arrow[rr,bend right]\&Y\arrow[r,hook]\&R^J
\end{tikzcd}}]}

\deleted{\uline{Prop} Soit $X,Y\in C_n$, alors $X\cap Y\in C_n$}

\uline{Prop} a) Si $X\in C_I$, $\operatorname{id}_X$ est $C_*$-morphisme

b) Le composé de deux $C_*$-morphismes est \unsure{itou}.

c) \deleted{Soit} $f:X\overset{\deleted{\unsure{\text{est}}}}{\longrightarrow}Y$ une bijection. Pour que $f$ soit un $C_*$-morphisme, il f. et s. que $f^{-1}$ le soit.

d) Soient $X\in C_I$, $Y\in C_J$, $Z\in C_K$, $g:X\to Y$, $h:X\to Z$, d'où $f=(g,h):X\to Y\times Z$ \deleted{$\unsure{\in C}$} (où $Y\times Z\in C_{J\amalg K}$).
Alors $f$ est un $C_*$-morphisme ssi $g$ et $h$ le sont.

\margin{e) Soit $\alpha:\deleted{\unsure{E}}\inserted{I}\to J$, d'où $\alpha^*:R^J\to R^I$, soit $X\in C_J$, d'où $Y=\alpha^*(X)\in C_I$. Alors $\alpha^*|X:X\to Y$ est un $C_*$-morphisme.
f) Si $X\subset Y$, $X,Y\in C_I$, alors l'inclusion $X\lhook\joinrel\longrightarrow Y$ est un $C_*$-morphisme.}

\margin{g) Si $f:X\to Y$ est un $C_*$-morphisme, $X\in C_n$, $Y\in C_m$, alors $f(X')\in C_m$ pour $X'\subset X$, $X'\in C_n$.
h) Si $f:X\to Y$ comme dessus, alors $f^{-1}(Y')\in C_n$ si $Y'\subset Y$, $Y'\in C_m$.}

\uline{Définition} \deleted{\unsure{Une}} Structure \inserted{$T$ de $C_*$-modèle} sur un ens \deleted{\unsure{E}} est la donnée, pour \deleted{tout} $n\in\mathbb N$ d'un ens $T_n$ d'applications injectives de $E$ dans $\mathbb R^n$, telle que
